\begin{frontmatter}

%% Title, authors and addresses

%% use the tnoteref command within \title for footnotes;
%% use the tnotetext command for the associated footnote;

%% use the fnref command within \author or \address for footnotes;
%% use the fntext command for the associated footnote;

%% use the corref command within \author or corresponding author footnotes;
%% use the cortext command for the associated footnote;
%% use the ead command for the email address,
%% and the form \ead[url] for the home page:
%%
%% \title{Title\tnoteref{label1}}
%% \tnotetext[label1]{}
%% \author{Name\corref{cor1}\fnref{label2}}
%% \ead{email address}
%% \ead[url]{home page}
%% \fntext[label2]{}
%% \cortext[cor1]{}
%% \address{Address\fnref{label3}}
%% \fntext[label3]{}

%\dochead{Article Header}
%% Use \dochead if there is an article header, e.g. \dochead{Short communication}

\title{Segmentation-Assisted Synthetic Augmentation for Oriented Vehicle Detection in Peruvian Traffic Surveillance}

%% use optional labels to link authors explicitly to addresses:
%% \author[label1,label2]{<author name>}
%% \address[label1]{<address>}
%% \address[label2]{<address>}

\author[First]{Alvaro Raul Quispe Condori\corref{cor1}}
\ead{aquispecondo@unsa.edu.pe}

\author[First]{Denise Andrea Huacani Jara}
\ead{dhuacanij@unsa.edu.pe}

\author[First]{Fabiana Francinet Pacheco Palo}
\ead{fpachecop@unsa.edu.pe}

\author[First]{Saul Andre Sivincha Machaca}
\ead{ssivincham@unsa.edu.pe}

\author[First]{Dolly Yadhira Mollo Chuquica\~na}
\ead{dmolloc@unsa.edu.pe}


%%COMMENTED\fntext[label2]{Footnote for author 1}
\cortext[cor1]{Corresponding author.}

%% ADDRESS
\address[First]{Universidad Nacional de San Agust{\'i}n, Arequipa, Per{\'u}.}

\begin{abstract}
Urban traffic cameras are valuable for mobility analysis, but vehicle detection in Lima traffic scenes is affected by dense flow, elevated viewpoints, shadows, occlusions, object orientation, and long-tailed class frequencies. This paper presents \textit{Segmentation-Assisted Synthetic Augmentation for Oriented Vehicle Detection in Peruvian Traffic Surveillance}, a data-centric ablation study that tests whether train-only, mask-guided copy-paste can add minority-class evidence without contaminating validation data. The workflow audits the source annotations, converts rotated boxes to YOLO-OBB labels, extracts real vehicle instances with Segment Anything Model masks, inserts them into safe training slots, trains YOLO26s-OBB conditions, and evaluates final artifacts with Macro AP-rIoU across nine classes and seven rotated-IoU thresholds after homography-based motion filtering. The augmentation release contains 1,501 synthetic image-label pairs and 3,391 inserted instances. In the completed final artifact set, all six conditions report 0.000 Macro AP-rIoU, 0.000 AP@0.50, and 0.000 AP@0.80, with zero true positives at rIoU 0.50. Improvement C reduces false positives among the synthetic/LaMa variants and reaches 98.9 FPS, while Improvement B shows strong internal YOLO validation mAP50 (0.929), but neither result becomes a final AP gain. The contribution is a reproducible OBB-preserving augmentation protocol and an auditable ablation framework that exposes the gap between training diagnostics and final rotated-box evaluation.
\end{abstract}

\begin{keyword}
%% keywords here, in the form: keyword \sep keyword

%% MSC codes here, in the form: \MSC code \sep code
%% or \MSC[2008] code \sep code (2000 is the default)

traffic surveillance \sep oriented object detection \sep segmentation-assisted augmentation \sep copy-paste augmentation
\end{keyword}

\end{frontmatter}
